\section{Narrative Text}

\subsection{Definition}
A narrative text is a text that tells a stroy involving characters and sequence of events. Narrative texts are generally written to entertain the reader, although they may also communicate a moral lesson or message.


Common examples include:


\begin{itemize}
\item Fairy tales
\item Legends
\item Fables
\item Myths
\item Folk tales
\item Fictional stories
\end{itemize}

\subsection{Purpose}

The primary purpose of a narrative text is to entertain or engage the reader through a story.


A narrative may also have a secondary purpose, such as communicating a moral lesson, explaining a cultural belief, or conveying a particular message.


For example, \textit{The Tortoise and the Hare} entertains the reader while also communicatinga lesson about perseverance and overconfidence.

\subsection{Generic Structure}

A narrative text commonly consists of six structural elements:
\begin{enumerate}
\item Orientation
\item Complication
\item Sequence of Events
\item Resolution
\item Coda
\end{enumerate}

The \textit{purpose} describes why the narrative is written, while the five elements above describe how the story is organized.

\subsubsection{Orientation}

The orientation introduces the basic circumstances of the story.


It commonly establishes:

\begin{itemize}
\item The characters
\item The setting
\item The time
\item The initial situations
\end{itemize}

For example:


\begin{quotation}
  Once upon a time, a hare and a tortoise lived in a forest.
\end{quotation}

This sentence introduces the main characters and the setting.

\subsubsection{Complication}

The complication introduces the main problem or conflict faced by the characters.


Without a complication, there is generally no central problem for the story to resolve.


For example:

\begin{quotation}
  The hare challenged the tortoise to a race because he believed that he was much faster.
\end{quotation}


The challenge creates the central conflict of the story.


\subsubsection{Sequence of Events}

The sequence of events describe what happens as the story develops.


Events are usually presented in chronological order.


For example:

\begin{enumerate}
\item The race begins.
\item The hare runs far ahead.
\item The hare becomes overconfident and takes a nap.
\item The tortoise continues moving.
\item The tortoise approaches the finish line.
\end{enumerate}

The sequence connects the complication to the resolution.

\subsubsection{Resolution}

The resolution presents the outcome of the main conflict.


For example:

\begin{quotation}
  The tortoise reaches the finish line before the hare and wins the race.
\end{quotation}

The resolution does not necessarily have to be happy. A narrative may have a positive, negative or unresolved outcome.

\subsubsection{Coda}

The coda is an optional closing section that states the lesson, moral, or message of the story.


For example:

\begin{quotation}
  The story teaches that consistent effort is more valuable than overconfidence.
\end{quotation}

Not every narrative contains an explicit coda. Sometimes the reader is expected to infer the message from the events and resolution.
